\documentclass[12pt,a4paper,oneside]{report}

%% ─── Encodage & langue ────────────────────────────────────────────────────────
\usepackage[T1]{fontenc}
\usepackage[utf8]{inputenc}
\usepackage[french]{babel}
\usepackage{lmodern}
\usepackage{microtype}
\usepackage{svg}

%% ─── Mise en page ─────────────────────────────────────────────────────────────
\usepackage{geometry}
\usepackage{setspace}
\usepackage{emptypage}
\geometry{top=2.5cm, bottom=2.5cm, left=3cm, right=2.5cm}
\onehalfspacing
\setlength{\headheight}{24pt}
\setlength{\emergencystretch}{3em}

%% ─── Tableaux ─────────────────────────────────────────────────────────────────
\usepackage{array}
\usepackage{longtable}
\usepackage{svg}
\usepackage{tabularx}
\usepackage{booktabs}
\usepackage{multirow}

%% ─── Couleurs & graphiques ────────────────────────────────────────────────────
\usepackage{xcolor}
\usepackage{colortbl}
\usepackage{graphicx}
\usepackage{float}
\usepackage{caption}
\captionsetup{labelfont=bf, textfont=it, font=small}

%% ─── Style typographique des titres ───────────────────────────────────────────
\usepackage{titlesec}
\usepackage{titletoc}

%% ─── En-têtes et pieds de page ────────────────────────────────────────────────
\usepackage{fancyhdr}

%% ─── Boîtes colorées ──────────────────────────────────────────────────────────
\usepackage{tcolorbox}
\tcbuselibrary{skins, breakable}

%% ─── Références & bibliographie ───────────────────────────────────────────────
\usepackage[hidelinks]{hyperref}
\usepackage{csquotes}
\usepackage[backend=bibtex, style=numeric, sorting=none]{biblatex}
\usepackage{filecontents}
\begin{filecontents*}{inline_references.bib}
@online{unesco_digital_2026,
  author   = {{UNESCO}},
  title    = {AI and technologies in education (Digital Education)},
  year     = {2026},
  url      = {https://www.unesco.org/en/digital-education},
  urldate  = {2026-05-01}
}

@online{unesco_gem_2023,
  author   = {{UNESCO}},
  title    = {Global Education Monitoring Report: Summary 2023},
  year     = {2023},
  url      = {https://www.unesco.org/gem-report/sites/default/files/medias/fichiers/2023/07/Summary_v5.pdf},
  urldate  = {2026-05-01}
}

@online{unesco_oer_2019,
  author   = {{UNESCO}},
  title    = {Recommendation on Open Educational Resources (OER)},
  year     = {2019},
  url      = {https://unesdoc.unesco.org/ark:/48223/pf0000373755},
  urldate  = {2026-05-01}
}

@online{itu_facts_2025,
  author   = {{International Telecommunication Union}},
  title    = {Measuring digital development: Facts and Figures 2025},
  year     = {2025},
  url      = {https://www.itu.int/en/ITU-D/Statistics/Pages/facts/default.aspx},
  urldate  = {2026-05-01}
}

@online{worldbank_edtech_2026,
  author   = {{World Bank}},
  title    = {Digital Technologies in Education},
  year     = {2026},
  url      = {https://www.worldbank.org/en/topic/edutech},
  urldate  = {2026-05-01}
}

@online{moodle_overview_2026,
  author   = {{Moodle}},
  title    = {Moodle: Learning Management System Overview},
  year     = {2026},
  url      = {https://moodle.com/},
  urldate  = {2026-05-01}
}

@online{google_classroom_2026,
  author   = {{Google for Education}},
  title    = {Google Classroom / Workspace for Education Overview},
  year     = {2026},
  url      = {https://edu.google.com/workspace-for-education/},
  urldate  = {2026-05-01}
}

@online{microsoft_teams_edu_2026,
  author   = {{Microsoft Education}},
  title    = {Microsoft Teams for Education Overview},
  year     = {2026},
  url      = {https://www.microsoft.com/education/products/teams},
  urldate  = {2026-05-01}
}

@online{canvas_overview_2026,
  author   = {{Instructure}},
  title    = {Canvas LMS Overview},
  year     = {2026},
  url      = {https://www.instructure.com/canvas},
  urldate  = {2026-05-01}
}

@misc{memoire_branding_2025,
  author       = {Ouldbrahem Rajaa},
  title        = {Une application mobile de Branding, Mémoire de fin d'études},
  year         = {2025},
  howpublished = {USTO-MB, document interne}
}

@misc{memoire_bibliotheque_2025,
  author       = {Bessahraoui Sarah and Neggal Aya},
  title        = {Prise en place d'une bibliothèque virtuelle pour le partage de livres, Mémoire de fin d'études},
  year         = {2025},
  howpublished = {USTO-MB, document interne}
}

@book{osterwalder2010bmc,
  author    = {Alexander Osterwalder and Yves Pigneur},
  title     = {Business Model Generation},
  year      = {2010},
  publisher = {John Wiley \& Sons}
}

@book{kotler2017marketing,
  author    = {Philip Kotler and Kevin Lane Keller},
  title     = {Marketing Management},
  edition   = {15},
  year      = {2017},
  publisher = {Pearson}
}

@book{ries2011lean,
  author    = {Eric Ries},
  title     = {The Lean Startup},
  year      = {2011},
  publisher = {Crown Business}
}
\end{filecontents*}
\addbibresource{inline_references.bib}

%% ═══════════════════════════════════════════════════════════════════════════════
%% PALETTE DE COULEURS
%% ═══════════════════════════════════════════════════════════════════════════════
\definecolor{ustoblue}{RGB}{0,51,102}
\definecolor{ustolightblue}{RGB}{210,228,248}
\definecolor{edugreen}{RGB}{195,235,200}
\definecolor{edudarkgreen}{RGB}{0,100,0}
\definecolor{rowgray}{RGB}{245,245,245}
\definecolor{headerwhite}{RGB}{255,255,255}

%% ═══════════════════════════════════════════════════════════════════════════════
%% STYLE DES CHAPITRES, SECTIONS, SOUS-SECTIONS
%% ═══════════════════════════════════════════════════════════════════════════════
\titleformat{\chapter}[block]
  {\normalfont\LARGE\bfseries\color{ustoblue}}
  {Chapitre~\thechapter~\textemdash~}{0pt}{\LARGE\bfseries}
\titlespacing*{\chapter}{0pt}{0pt}{28pt}

\titleformat{\section}
  {\normalfont\large\bfseries\color{ustoblue}}
  {\thesection}{1em}{}
\titlespacing*{\section}{0pt}{14pt}{6pt}

\titleformat{\subsection}
  {\normalfont\normalsize\bfseries\color{black}}
  {\thesubsection}{1em}{}
\titlespacing*{\subsection}{0pt}{10pt}{4pt}

\titleformat{\paragraph}[runin]
  {\normalfont\normalsize\bfseries}
  {}{0pt}{}[~\textemdash~]

%% ═══════════════════════════════════════════════════════════════════════════════
%% EN-TÊTES ET PIEDS DE PAGE
%% ═══════════════════════════════════════════════════════════════════════════════
\pagestyle{fancy}
\fancyhf{}
\fancyhead[L]{\parbox[c]{0.44\headwidth}{\small\itshape USTO-MB~---~D\'ep. Informatique}}
\fancyhead[R]{\parbox[c]{0.44\headwidth}{\small\itshape\raggedleft\nouppercase{\leftmark}}}
\fancyfoot[C]{\small ---~\thepage~---}
\renewcommand{\headrulewidth}{0.5pt}
\renewcommand{\footrulewidth}{0.3pt}

% Style pour pages spéciales (TOC, LOF, LOT)
\fancypagestyle{plain}{%
  \fancyhf{}%
  \fancyfoot[C]{\small ---~\thepage~---}%
  \renewcommand{\headrulewidth}{0pt}%
}

%% ═══════════════════════════════════════════════════════════════════════════════
%% BOÎTE D'INSTRUCTIONS FIGURE (à supprimer avant soumission finale)
%% ═══════════════════════════════════════════════════════════════════════════════
\newtcolorbox{figurebox}[1]{
  colback  = ustolightblue!35,
  colframe = ustoblue,
  coltitle = white,
  fonttitle= \bfseries\small,
  title    = {Instruction figure~---~#1},
  breakable,
  left=6pt, right=6pt, top=4pt, bottom=4pt,
  before=\medskip, after=\medskip,
  sharp corners=south
}

%% ═══════════════════════════════════════════════════════════════════════════════
%% COMMANDE IMAGE SÉCURISÉE (placeholder si fichier absent)
%% ═══════════════════════════════════════════════════════════════════════════════
\newcommand{\safeimage}[4][0.85\textwidth]{%
  \begin{figure}[H]
    \centering
    \IfFileExists{#2}{%
      \includegraphics[width=#1]{#2}%
    }{%
      \fcolorbox{ustoblue}{ustolightblue!20}{%
        \parbox[c][3.5cm][c]{0.88\textwidth}{\centering
          \color{ustoblue}\textbf{[Emplacement r\'eserv\'e~---~Figure~#3]}\\[6pt]
          \small Fichier attendu~:~\texttt{\detokenize{#2}}\\[4pt]
          \textit{\color{gray}Remplacer par l'image avant soumission}
        }%
      }%
    }%
    \caption{#4}
    \label{fig:#3}
  \end{figure}%
}

%% ═══════════════════════════════════════════════════════════════════════════════
%% DOCUMENT
%% ═══════════════════════════════════════════════════════════════════════════════
\begin{document}

% === Dédicaces et Remerciements (à personnaliser) ===
\cleardoublepage
\thispagestyle{empty}\chapter*{Dédicaces}
\addcontentsline{toc}{chapter}{Dédicaces}
\vspace*{3.5cm}
\begin{center}
\large
\textit{\`A ma famille, pour leur amour inconditionnel, leur patience et leur soutien constant tout au long de mon parcours.}
\medskip

\textit{\`A mes amis et coll\`egues, pour leur pr\'esence, leurs encouragements sinc\`eres et la richesse de nos \'echanges.}
\medskip

\textit{\`A toutes les personnes qui ont partag\'e leurs connaissances, leur temps et leur bienveillance pour m'aider \`a r\'ealiser ce travail.}
\end{center}
\vspace{1cm}
\clearpage

\thispagestyle{empty}
\chapter*{Remerciements}
\addcontentsline{toc}{chapter}{Remerciements}
\vspace*{3.5cm}
\begin{center}
\large
    \textit{Je tiens \`a exprimer ma profonde gratitude \`a mon encadrant pour la qualit\'e de son suivi, la pertinence de ses conseils ainsi que sa rigueur scientifique tout au long de ce travail.}
    
    \medskip
    \textit{Mes remerciements s'adressent aussi \`a l'ensemble de l'\'équipe p\'edagogique et au personnel administratif pour leur disponibilit\'e et leur soutien constant. }
    
    \medskip
    \textit{Enfin, j'adresse ma reconnaissance \`a mes coll\`egues et \`a ma famille pour leur soutien moral, leur patience et leurs encouragements pr\'ecieux tout au long de ce parcours.}
\end{center}
\clearpage
\pagenumbering{roman}
\tableofcontents
\clearpage
\listoffigures
\clearpage
\listoftables
\clearpage

\pagenumbering{arabic}
\setcounter{page}{1}

%% ═══════════════════════════════════════════════════════════════════════════════
%% INTRODUCTION GÉNÉRALE
%% ═══════════════════════════════════════════════════════════════════════════════

\chapter*{Introduction générale}
\addcontentsline{toc}{chapter}{Introduction générale}

La transformation numérique des systèmes d'enseignement constitue aujourd'hui une priorité stratégique pour les établissements confrontés à la massification des effectifs, à la diversification des parcours et à l'exigence croissante de qualité pédagogique. Dans ce contexte, la mise en place de plateformes numériques ne peut plus se limiter à la simple diffusion de supports de cours: elle doit intégrer la gouvernance académique, la collaboration entre acteurs et la traçabilité des activités.

Le projet \textbf{EduPlatform} s'inscrit dans cette dynamique. Il vise à proposer une plateforme unifiée de partage pédagogique et de collaboration, adaptée aux établissements scolaires, aux écoles de soutien et aux écoles de langues, en combinant: gestion par rôles, structuration des ressources par niveau/module/composant, communication synchrone et asynchrone, et outils de pilotage administratifs.

\paragraph{Problématique générale.}
Comment concevoir une plateforme éducative à la fois robuste techniquement, pertinente pédagogiquement et viable économiquement, capable d'unifier les usages quotidiens des apprenants, formateurs et responsables d'établissement?

\paragraph{Objectif global.}
Concevoir, implémenter et évaluer EduPlatform comme solution complète pour les écoles et centres de formation, puis proposer un cadre de passage à l'échelle startup à travers un modèle économique et une stratégie de croissance.

\paragraph{Organisation du mémoire.}
Ce mémoire est structure en quatre chapitres:
\begin{itemize}
  \item \textbf{Chapitre 1}: Etude de contexte, problématique et analyse des besoins;
  \item \textbf{Chapitre 2}: Conception UML et architecture de la solution;
  \item \textbf{Chapitre 3}: Implementation technique, deploiement et validation;
  \item \textbf{Chapitre 4}: Evaluation stratégique et business model (cadre startup, loi 12-75).
\end{itemize}

L'enchainement de ces chapitres permet de relier de maniere progressive la pertinence scientifique du projet, sa realisation concrete et sa projection entrepreneuriale.


%% ═══════════════════════════════════════════════════════════════════════════════
%% CHAPITRE 1 — Contexte général, problématique et état de l'art
%% ═══════════════════════════════════════════════════════════════════════════════

\chapter{Etude de contexte, problématique et analyse des besoins}

%%──────────────────────────────────────────────────────────────────────────────
\section{Introduction du chapitre}
%%──────────────────────────────────────────────────────────────────────────────

La g\'en\'eralisation des technologies num\'eriques dans l'enseignement et la formation constitue aujourd'hui l'un des leviers strat\'egiques les plus significatifs pour l'am\'elioration de la qualit\'e p\'edagogique, la continuit\'e des activit\'es et l'inclusion \'educative. Dans ce contexte de mutation profonde, les \'etablissements scolaires, les centres de soutien et les écoles de langues sont confront\'es \`a un imp\'eratif d'adaptation organisationnelle et technologique, auquel r\'epondent des environnements num\'eriques de plus en plus sophistiqu\'es~\cite{unesco_digital_2026,worldbank_edtech_2026}.

Le pr\'esent chapitre constitue le socle analytique et conceptuel du projet \textbf{EduPlatform}. Il poursuit quatre objectifs compl\'ementaires~: (i)~d\'ecrire le contexte mondial et national de la transformation num\'erique de l'enseignement~; (ii)~formuler la probl\'ematique ayant motiv\'e la conception d'EduPlatform~; (iii)~pr\'esenter les objectifs et la m\'ethodologie adopt\'es~; et (iv)~dresser un \'etat de l'art approfondi des solutions existantes, afin de positionner EduPlatform par rapport aux r\'ef\'erences du domaine.

Ce chapitre est organis\'e comme suit~: la section~\ref{sec:contexte} traite du contexte g\'en\'eral~; la section~\ref{sec:problématique} formule la probl\'ematique~; la section~\ref{sec:objectifs} d\'etaille les objectifs~; la section~\ref{sec:methodo} pr\'esente la d\'emarche m\'ethodologique~; la section~\ref{sec:etat_art} constitue l'\'etat de l'art~; la section~\ref{sec:positionnement} positionne EduPlatform~; la section~\ref{sec:figures} pr\'esente les figures requises~; enfin, la section~\ref{sec:conclusion} conclut le chapitre.

%%──────────────────────────────────────────────────────────────────────────────
\section{Contexte g\'en\'eral}
\label{sec:contexte}
%%──────────────────────────────────────────────────────────────────────────────

\subsection{Transformation num\'erique de l'enseignement sup\'erieur \`a l'\'echelle mondiale}

La transformation num\'erique de l'\'education s'est acc\'el\'er\'ee de mani\`ere consid\'erable au cours de la derni\`ere d\'ecennie, port\'ee par la g\'en\'eralisation de l'acc\`es \`a Internet, la mont\'ee en puissance des technologies mobiles et, plus r\'ecemment, l'\'emergence de l'intelligence artificielle appliqu\'ee aux processus p\'edagogiques. L'Union internationale des t\'el\'ecommunications (UIT) indique que 6~milliards de personnes, soit environ 74\,\% de la population mondiale, utilisaient Internet en 2025, tout en soulignant que 2,2~milliards d'individus demeurent exclus de cet acc\`es, principalement dans les pays \`a revenus faibles et interm\'ediaires~\cite{itu_facts_2025}.

L'UNESCO confirme que si la connectivit\'e progresse \`a l'\'echelle mondiale, elle reste profond\'ement in\'egale~: seulement 40\,\% des \'écoles primaires et 50\,\% des \'etablissements du secondaire inf\'erieur sont connect\'es \`a Internet dans le monde~\cite{unesco_gem_2023}. Ces donn\'ees soul\`event des questions fondamentales de politique \'educative et d'\'equit\'e d'acc\`es aux ressources num\'eriques, que l'UNESCO aborde dans ses cadres de comp\'etences num\'eriques et ses recommandations sur les ressources \'educatives libres~\cite{unesco_oer_2019,unesco_digital_2026}.

La Banque mondiale souligne que l'EdTech repr\'esente une opportunit\'e in\'edite pour les pays \`a revenus interm\'ediaires~: si elle est soutenue par des comp\'etences num\'eriques fondamentales et une connectivit\'e ad\'equate, elle peut contribuer \`a r\'eduire la crise d'apprentissage mondiale, notamment dans les communaut\'es souffrant de p\'enurie d'enseignants~\cite{worldbank_edtech_2026}. Toutefois, cette m\^eme institution met en garde contre les risques li\'es \`a la s\'ecurit\'e des donn\'ees et le risque d'aggravation des in\'egalit\'es \'educatives en l'absence de politiques d'encadrement adapt\'ees.

\subsection{Contexte national et enjeux institutionnels}

Dans le contexte alg\'erien et plus particuli\`erement au sein des écoles de soutien, des centres de formation et des écoles de langues, la gestion des activit\'es p\'edagogiques souffre d'une fragmentation des outils et des canaux de diffusion. Les ressources sont souvent partag\'ees via des messageries instantan\'ees non structur\'ees, les annonces ne parviennent pas syst\'ematiquement aux apprenants concern\'es, et les enseignants manquent d'un espace centralis\'e pour organiser leurs productions p\'edagogiques par module et par niveau d'\'études.

%%──────────────────────────────────────────────────────────────────────────────
\section{Probl\'ematique}
\label{sec:problématique}
%%──────────────────────────────────────────────────────────────────────────────

La gestion des activit\'es p\'edagogiques dans un contexte éducatif contemporain fait face \`a un d\'efi central~: la multiplicit\'e des outils et des canaux de communication g\'en\`ere une dispersion informationnelle qui nuit \`a la coh\'erence de l'exp\'erience des apprenants et des formateurs. Les ressources p\'edagogiques~--- cours, travaux dirig\'es, travaux pratiques, examens, corrig\'es et supports multim\'edias~--- sont dispers\'ees sur des supports h\'et\'erog\`enes, sans structuration par module ni tra\c{c}abilit\'e des mises \`a jour.

Par ailleurs, la communication entre enseignants, mod\'erateurs et apprenants s'appuie souvent sur des outils grand public non con\c{c}us pour les usages éducatifs, ce qui engendre des probl\`emes de s\'ecurit\'e, de gouvernance et de qualit\'e de service. L'absence d'un syst\`eme de notification centralis\'e, d'un espace de publication structur\'e et d'un planificateur int\'egr\'e contraint les acteurs \`a recourir \`a des solutions de contournement chronophages et peu fiables.

\medskip
\noindent\fcolorbox{ustoblue}{ustolightblue!20}{\parbox{0.96\textwidth}{%
\textbf{Question de recherche~:} Comment concevoir et d\'evelopper une plateforme num\'erique unifi\'ee, s\'ecuris\'ee et modulaire, capable de centraliser les ressources p\'edagogiques, de faciliter la communication entre les diff\'erents acteurs de la communaut\'e éducative, et de soutenir l'organisation du travail collectif, tout en restant accessible, maintenable et adapt\'ee au contexte éducatif local (écoles, centres de soutien, écoles de langues)?
}}
\medskip

%%──────────────────────────────────────────────────────────────────────────────
\section{Objectifs du projet}
\label{sec:objectifs}
%%──────────────────────────────────────────────────────────────────────────────

\subsection{Objectif g\'en\'eral}

L'objectif g\'en\'eral du projet est de concevoir, d\'evelopper et valider \textbf{EduPlatform}~: une plateforme web de partage p\'edagogique et de collaboration éducative, con\c{c}ue pour les usages des écoles, des centres de soutien et des écoles de langues, et structur\'ee autour d'une architecture orient\'ee modules, d'un syst\`eme de gouvernance par r\^oles et d'une couche sociale int\'egr\'ee.

\subsection{Objectifs sp\'ecifiques}

\begin{itemize}
  \item[\textbf{OS1.}] Centraliser les ressources p\'edagogiques (cours, TD, TP, examens, corrig\'es, playlists) dans une architecture hi\'erarchique~: Niveau $\rightarrow$ Module $\rightarrow$ Section $\rightarrow$ Ressource.
  \item[\textbf{OS2.}] Mettre en \oe uvre un syst\`eme d'authentification s\'ecuris\'e par jetons JWT et une autorisation par r\^oles~: Administrateur global, Mod\'erateur et Utilisateur.
  \item[\textbf{OS3.}] Int\'egrer un espace de communication structur\'e~: publications, commentaires, votes, partages et messagerie instantan\'ee (chat priv\'e et de groupe).
  \item[\textbf{OS4.}] D\'evelopper un syst\`eme de notifications permettant le suivi des \'ev\'enements acad\'emiques avec compteur de non-lus.
  \item[\textbf{OS5.}] Garantir la tra\c{c}abilit\'e compl\`ete des actions de gestion~: ajout, modification et suppression des ressources, contr\^ol\'ees par le syst\`eme de r\^oles.
\end{itemize}

%%──────────────────────────────────────────────────────────────────────────────
\section{D\'emarche m\'ethodologique}
\label{sec:methodo}
%%──────────────────────────────────────────────────────────────────────────────

La d\'emarche adopt\'ee suit un processus incr\'emental en cinq phases, conforme aux pratiques reconnues en ing\'enierie des syst\`emes d'information et aux recommandations des travaux de r\'ef\'erence analys\'es~:

\begin{enumerate}
  \item \textbf{Analyse des besoins}~: identification des besoins fonctionnels et non fonctionnels par l'\'étude des usages existants, l'analyse de la litt\'erature et la consultation des trois m\'emoires de r\'ef\'erence.
  \item \textbf{Mod\'elisation conceptuelle}~: repr\'esentation des interactions par des diagrammes UML (cas d'utilisation, s\'equence, classes) pour formaliser les sc\'enarios et l'architecture logique.
  \item \textbf{Conception de l'architecture}~: d\'efinition de l'architecture applicative (API REST, frontend SPA) et du sch\'ema de la base de donn\'ees relationnelle PostgreSQL.
  \item \textbf{Impl\'ementation et int\'egration}~: d\'eveloppement progressif par fonctionnalit\'es (authentification, ressources, social, notifications, planificateur, administration), avec validation continue.
  \item \textbf{Tests et ajustements}~: validation fonctionnelle, tests de s\'ecurit\'e et ajustements selon les retours d'exp\'erience.
\end{enumerate}

%%──────────────────────────────────────────────────────────────────────────────
\section{\'Etat de l'art}
\label{sec:etat_art}
%%──────────────────────────────────────────────────────────────────────────────

\subsection{D\'efinitions et concepts cl\'es}

\paragraph{Syst\`eme de gestion de l'apprentissage (LMS).}
Un LMS (\textit{Learning Management System}) est un logiciel permettant de cr\'eer, distribuer et g\'erer du contenu p\'edagogique, ainsi que de suivre la progression des apprenants~\cite{moodle_overview_2026}. Les principaux LMS du march\'e incluent Moodle, Canvas et Blackboard.

\paragraph{Plateforme collaborative.}
Une plateforme collaborative est un environnement num\'erique qui offre des fonctionnalit\'es de communication, de co-production et de gestion des activit\'es pour des groupes d'utilisateurs. Elle peut int\'egrer des outils de messagerie, de partage de fichiers, de planification et de gestion de projets.

\paragraph{Apprentissage en ligne (\textit{e-learning}).}
L'e-learning d\'esigne l'ensemble des modalit\'es d'enseignement et d'apprentissage faisant appel aux technologies num\'eriques, qu'elles soient synchrones (classes virtuelles, visioconf\'erence) ou asynchrones (cours en ligne, forums, documents partag\'es)~\cite{unesco_digital_2026}.

\subsection{Pr\'esentation des solutions de r\'ef\'erence}

\paragraph{Moodle.}
Moodle est l'un des LMS open source les plus utilis\'es dans le monde. Il offre une gestion fine des cours, des activit\'es et des \'évaluations, ainsi qu'un syst\`eme d'extensions (plugins) permettant une forte personnalisation. Sa gouvernance institutionnelle est tr\`es flexible, mais son d\'eploiement exige des comp\'etences techniques~\cite{moodle_overview_2026}.

\paragraph{Google Classroom.}
Google Classroom est un service cloud de gestion de classes orient\'e simplicit\'e. Il s'int\`egre nativement aux outils Google Workspace (Drive, Docs, Meet). Sa personnalisation reste cependant limit\'ee et sa d\'ependance \`a l'\'ecosyst\`eme propri\'etaire peut constituer une contrainte~\cite{google_classroom_2026}.

\paragraph{Microsoft Teams pour l'\'Education.}
Microsoft Teams Education offre une exp\'erience unifi\'ee de classe virtuelle et de collaboration au sein de l'environnement Microsoft~365. Il supporte la visioconf\'erence, le chat, la gestion d'\'équipes et les devoirs~\cite{microsoft_teams_edu_2026}.

\paragraph{Canvas LMS.}
Canvas est un LMS cloud moderne, appr\'eci\'e pour son interface structur\'ee, ses workflows d'\'évaluation et ses outils analytiques. Son adoption est pr\'edominante dans les établissements d'enseignement nord-am\'ericains~\cite{canvas_overview_2026}.

%%──────────────────────────────────────────────────────────────────────────────
\subsection{Tableau~1~---~Comparaison synth\'etique des plateformes}
%%──────────────────────────────────────────────────────────────────────────────

Le tableau~\ref{tab:compar_synth} offre une vue d'ensemble des quatre plateformes de r\'ef\'erence selon quatre axes strat\'egiques~: positionnement, points forts, limites et pertinence dans le contexte éducatif local algérien.

\renewcommand{\arraystretch}{1.8}
{\setlength{\tabcolsep}{7pt}
\begin{table}[H]
\centering
\caption{Tableau synth\'etique des plateformes p\'edagogiques num\'eriques (vue d'ensemble)}
\label{tab:compar_synth}
\begin{tabular}{|>{\bfseries}p{2.3cm}|p{2.5cm}|p{3.2cm}|p{3.0cm}|p{2.2cm}|}

%% ── EN-TÊTE ────────────────────────────────────────────────────────────
\hline
\rowcolor{ustoblue}
\textcolor{white}{\textbf{Plateforme}} &
\textcolor{white}{\textbf{Positionnement}} &
\textcolor{white}{\textbf{Points forts}} &
\textcolor{white}{\textbf{Limites}} &
\textcolor{white}{\textbf{Pertinence locale}} \\
\hline

%% ── LIGNE 1 ────────────────────────────────────────────────────────────
\rowcolor{ustolightblue!40}
Moodle &
LMS open source &
Flexibilit\'e, plugins, structure p\'edagogique fine &
Maintenance technique exigeante &
\'Elev\'ee (mais complexe) \\
\hline

%% ── LIGNE 2 ────────────────────────────────────────────────────────────
\rowcolor{rowgray}
Google Classroom &
Classe et devoirs cloud &
Simplicit\'e, int\'egration Google Workspace &
D\'ependance Google, personnalisation limit\'ee &
Moyenne \\
\hline

%% ── LIGNE 3 ────────────────────────────────────────────────────────────
\rowcolor{ustolightblue!40}
MS Teams Edu &
Collaboration + classe &
Communication unifi\'ee, visioconf\'erence &
Complexit\'e pour usages purement p\'edagogiques &
Moyenne \\
\hline

%% ── LIGNE 4 ────────────────────────────────────────────────────────────
\rowcolor{rowgray}
Canvas LMS &
LMS cloud moderne &
Interface structur\'ee, analytics &
Co\^ut selon contexte institutionnel &
Faible \`a moyenne \\
\hline

\end{tabular}
\end{table}}

\noindent{\small\textit{Sources~:~\cite{moodle_overview_2026,google_classroom_2026,microsoft_teams_edu_2026,canvas_overview_2026}}}

%%──────────────────────────────────────────────────────────────────────────────
%% BLOC DE TRANSITION ENTRE LES DEUX TABLEAUX
%%──────────────────────────────────────────────────────────────────────────────

\medskip
\begin{tcolorbox}[
  colback  = ustolightblue!20,
  colframe = ustoblue,
  coltitle = white,
  fonttitle= \bfseries,
  title    = {Observations cl\'es issues du tableau~\ref{tab:compar_synth}~---~vers une analyse crit\`ere par crit\`ere},
  left=8pt, right=8pt, top=6pt, bottom=6pt,
  sharp corners=south
]
\begin{itemize}
  \item Aucune des quatre solutions g\'en\'eralistes n'est con\c{c}ue sp\'ecifiquement pour le contexte éducatif algérien~; leur adaptation exige des ressources techniques et financi\`eres importantes.
  \item EduPlatform est la seule solution \textbf{auto-h\'eberg\'ee} et \textbf{contexte-native}~: chaque fonctionnalit\'e a \'et\'e con\c{c}ue en r\'eponse directe \`a un besoin identifi\'e dans un contexte éducatif local.
  \item Le tableau~\ref{tab:compar_detail} ci-dessous approfondit l'analyse sur \textbf{12 crit\`eres fonctionnels et techniques}, permettant de mesurer pr\'ecis\'ement les apports sp\'ecifiques d'EduPlatform.
\end{itemize}
\end{tcolorbox}
\medskip

%%──────────────────────────────────────────────────────────────────────────────
\subsection{Tableau~2~---~Analyse comparative d\'etaill\'ee (12 crit\`eres)}
%%──────────────────────────────────────────────────────────────────────────────

Le tableau~\ref{tab:compar_detail} confronte les quatre plateformes de r\'ef\'erence sur douze crit\`eres couvrant les dimensions fonctionnelles, techniques, organisationnelles et \'economiques.

\renewcommand{\arraystretch}{1.8}
\setlength{\tabcolsep}{7pt}
\begin{longtable}{|>{\bfseries\raggedright\arraybackslash}p{3.5cm}|
                   >{\raggedright\arraybackslash}p{2.5cm}|
                   >{\raggedright\arraybackslash}p{2.5cm}|
                   >{\raggedright\arraybackslash}p{2.5cm}|
                   >{\raggedright\arraybackslash}p{2.5cm}|}

\caption{Analyse comparative d\'etaill\'ee des plateformes p\'edagogiques num\'eriques (12 crit\`eres)}
\label{tab:compar_detail} \\

%% ── EN-TÊTE PREMIÈRE PAGE ──────────────────────────────────────────────
\hline
\rowcolor{ustoblue}
\textcolor{white}{\textbf{Crit\`ere}} &
\textcolor{white}{\textbf{Moodle}} &
\textcolor{white}{\textbf{G. Classroom}} &
\textcolor{white}{\textbf{MS Teams}} &
\textcolor{white}{\textbf{Canvas}} \\
\hline
\endfirsthead

%% ── EN-TÊTE PAGES SUIVANTES ────────────────────────────────────────────
\multicolumn{5}{l}{\small\textit{(suite du tableau~\ref{tab:compar_detail})}} \\[2pt]
\hline
\rowcolor{ustoblue}
\textcolor{white}{\textbf{Crit\`ere}} &
\textcolor{white}{\textbf{Moodle}} &
\textcolor{white}{\textbf{G. Classroom}} &
\textcolor{white}{\textbf{MS Teams}} &
\textcolor{white}{\textbf{Canvas}} \\
\hline
\endhead

\hline
\multicolumn{5}{r|}{\small\textit{Suite page suivante\ldots}} \\
\endfoot

\hline
\endlastfoot

%% ── CRITÈRE 1 ──────────────────────────────────────────────────────────
\rowcolor{ustolightblue!40}
Mode de d\'eploiement &
Auto-h\'eberg\'e\,/ cloud &
Cloud Google &
Cloud Microsoft &
Cloud \\
\hline

%% ── CRITÈRE 2 ──────────────────────────────────────────────────────────
\rowcolor{rowgray}
Orientation principale &
LMS complet &
Classe et devoirs &
Collaboration + classe &
LMS complet \\
\hline

%% ── CRITÈRE 3 ──────────────────────────────────────────────────────────
\rowcolor{ustolightblue!40}
Gestion des ressources &
Avanc\'ee (cours, banques) &
Essentielle (devoirs) &
Correcte (canaux) &
Avanc\'ee \\
\hline

%% ── CRITÈRE 4 ──────────────────────────────────────────────────────────
\rowcolor{rowgray}
Organisation par modules &
Partielle (cat\'eg.) &
Par classe uniquement &
Par \'équipe\,/\,canal &
Par cours \\
\hline

%% ── CRITÈRE 5 ──────────────────────────────────────────────────────────
\rowcolor{ustolightblue!40}
Syst\`eme de r\^oles &
Oui (multiples) &
Enseignant\,/\, \'El\`eve &
Propri\'etaire\,/\, Membre &
Admin\,/\, \'Etudiant \\
\hline

%% ── CRITÈRE 6 ──────────────────────────────────────────────────────────
\rowcolor{rowgray}
Publications et commentaires &
Forums de discussion &
Commentaires de flux &
Canaux de discussion &
Annonces de cours \\
\hline

%% ── CRITÈRE 7 ──────────────────────────────────────────────────────────
\rowcolor{ustolightblue!40}
Chat int\'egr\'e &
Via plugin &
Non (externe) &
Oui (complet) &
Partiel \\
\hline

%% ── CRITÈRE 8 ──────────────────────────────────────────────────────────
\rowcolor{rowgray}
Notifications &
Courriel\,/\, interne &
Courriel Google &
Notifications M365 &
Courriel\,/\, interne \\
\hline

%% ── CRITÈRE 9 ──────────────────────────────────────────────────────────
\rowcolor{ustolightblue!40}
Planificateur &
Agenda basique &
Google Agenda (ext.) &
Calendrier M365 &
Calendrier int\'egr\'e \\
\hline

%% ── CRITÈRE 10 ─────────────────────────────────────────────────────────
\rowcolor{rowgray}
Technologies &
PHP, MySQL &
Cloud Google &
Cloud Microsoft &
Ruby, PostgreSQL \\
\hline

%% ── CRITÈRE 11 ─────────────────────────────────────────────────────────
\rowcolor{ustolightblue!40}
Co\^ut\,/\,Licence &
Gratuit (open source) &
Gratuit (Google) &
Payant (M365) &
Payant (licence) \\
\hline

%% ── CRITÈRE 12 ─────────────────────────────────────────────────────────
\rowcolor{rowgray}
Adaptation locale &
\'Elev\'ee avec configuration &
Moyenne &
Moyenne &
Faible \\
\hline

\end{longtable}

\noindent{\small\textit{Sources~: \cite{moodle_overview_2026,google_classroom_2026,microsoft_teams_edu_2026,canvas_overview_2026}}}

\subsection{Synth\`ese et bilan de l'\'etat de l'art}

L'analyse comparative men\'ee dans les tableaux~\ref{tab:compar_synth} et~\ref{tab:compar_detail} confirme que les plateformes g\'en\'eralistes existantes couvrent partiellement les besoins identifi\'es dans notre contexte~: Moodle excelle dans la gestion p\'edagogique mais exige une expertise technique importante~; Google Classroom privil\'egie la simplicit\'e au d\'etriment de la personnalisation~; Microsoft Teams est optimis\'e pour la collaboration mais reste secondaire en mati\`ere d'organisation p\'edagogique fine~; enfin, Canvas est trop co\^uteux et peu adapt\'e aux contraintes institutionnelles locales.

Ces lacunes cumul\'ees justifient la conception d'EduPlatform~: une solution \'elabor\'ee sp\'ecifiquement pour le contexte acad\'emique local, gratuite, auto-h\'eberg\'ee, et combinant dans un seul environnement la gestion modulaire des ressources, la gouvernance par r\^oles, la communication int\'egr\'ee et la planification hebdomadaire.

%%──────────────────────────────────────────────────────────────────────────────
\section{Positionnement d'EduPlatform}
\label{sec:positionnement}
%%──────────────────────────────────────────────────────────────────────────────

EduPlatform se positionne comme une plateforme acad\'emique int\'egr\'ee, r\'eunissant dans un seul environnement les fonctions p\'edagogiques, sociales et organisationnelles que les solutions existantes dispersent entre plusieurs outils. Son architecture repose sur quatre piliers structurants~:

\begin{itemize}
  \item \textbf{Navigation module-first~:} la granularit\'e p\'edagogique est au c\oe ur de l'exp\'erience~; chaque module dispose d'une page d\'edi\'ee regroupant ses sections et ses ressources class\'ees par type.
  \item \textbf{Gouvernance par r\^oles~:} le syst\`eme d'autorisation \`a cinq niveaux (Administrateur Super, Administrateur Ecole, Formateur Principal, Formateur Simple, Utilisateur) garantit la s\'ecurit\'e, la modulation des acc\`es et la responsabilit\'e de chaque acteur.
  \item \textbf{Couche sociale et collaborative~:} le chat et les notifications constituent un tissu de communication acad\'emique structur\'e, tra\c{c}able et centralis\'e.
\end{itemize}

%%──────────────────────────────────────────────────────────────────────────────
\section{Illustrations de la plateforme EduPlatform}
\label{sec:figures}
%%──────────────────────────────────────────────────────────────────────────────

Cette section pr\'esente les captures d'\'ecran de la plateforme illustrant les fonctionnalit\'es d\'ecrites dans les sections pr\'ec\'edentes.

%%──────────────────────────
\subsection{Figure~1~---~Capture~: navigation par modules}

\safeimage[0.85\textwidth]{figures/F1_interface_modules.png}{F1}{Interface principale d'EduPlatform~: navigation par modules (vue liste des modules par niveau).}

%%──────────────────────────
\subsection{Figure~2~---~Capture~: ressources dans un module}



\safeimage[0.85\textwidth]{figures/F2_interface_ressources.png}{F2}{Interface d'EduPlatform~: page d\'edi\'ee d'un module avec ses ressources organis\'ees par section (vue administrateur).}

%%──────────────────────────
\subsection{Figure~3~---~Capture~: publications et chat}


\safeimage[0.85\textwidth]{figures/F3_interface_social.png}{F3}{Interface d'EduPlatform~: fil de publications et chat int\'egr\'es~--- couche sociale de la plateforme.}

%%──────────────────────────────────────────────────────────────────────────────
\section{Conclusion}
\label{sec:conclusion}
%%──────────────────────────────────────────────────────────────────────────────

Ce premier chapitre a \'etabli le cadre analytique et conceptuel du projet EduPlatform en articulant trois niveaux de lecture compl\'ementaires~: le contexte mondial de la transformation num\'erique de l'enseignement, les besoins sp\'ecifiques du contexte acad\'emique local de l'USTO-MB, et le positionnement d'EduPlatform par rapport aux solutions existantes.

L'analyse comparative men\'ee dans la section~\ref{sec:etat_art} d\'emontre qu'aucune des plateformes \'etudi\'ees~--- Moodle, Google Classroom, Microsoft Teams Education, Canvas~--- ne couvre de mani\`ere optimale l'ensemble des exigences identifi\'ees~: navigation module-first, gouvernance \`a trois r\^oles, communication int\'egr\'ee et planification hebdomadaire d\'edi\'ee. Cette convergence de lacunes justifie pleinement la conception d'EduPlatform comme solution \textit{ad hoc}, d\'evelopp\'ee sp\'ecifiquement pour le contexte éducatif local.

La suite du m\'emoire approfondit successivement la conception UML (Chapitre~2), l'impl\'ementation technique (Chapitre~3) et l'\'évaluation des r\'esultats (Chapitre~4), en s'appuyant sur le cadre th\'eorique et m\'ethodologique pos\'e dans le pr\'esent chapitre. Cette analyse vise notamment les besoins des écoles de soutien, des écoles de langues et des centres de formation.

%% ═══════════════════════════════════════════════════════════════════════════════
%% CHAPITRE 2 — Analyse, modelisation UML et conception
%% ═══════════════════════════════════════════════════════════════════════════════

\chapter{Conception UML et architecture de la solution}

%%──────────────────────────────────────────────────────────────────────────────
\section{Introduction du chapitre}
%%──────────────────────────────────────────────────────────────────────────────

Ce chapitre formalise la traduction de la problématique metier en specifications exploitables pour la realisation logicielle. Il presente successivement l'analyse des besoins, la modelisation UML des interactions, la conception des données et les choix architecturaux retenus pour EduPlatform.

L'objectif est double: garantir la cohérence fonctionnelle de la solution et réduire les risques de derivation durant l'implementation.

%%──────────────────────────────────────────────────────────────────────────────
\section{Analyse des besoins}
%%──────────────────────────────────────────────────────────────────────────────

\subsection{Identification des acteurs}

Quatre acteurs principaux structurent le système:
\begin{itemize}
  \item \textbf{Super administrateur}: gouvernance plateforme, écoles, administrateurs;
  \item \textbf{Administrateur d'école}: gestion des utilisateurs, modules, assignations, inscriptions;
  \item \textbf{Formateur Globale}: publication des ressources, suivi des apprenants, sessions live;
  \item \textbf{Formateur Simple}: publication des ressources, suivi des apprenants, sessions live;
  \item \textbf{Etudiant}: consultation des ressources, progression, interactions et lives.
\end{itemize}

\subsection{Besoins fonctionnels majeurs}

\begin{itemize}
  \item authentification securisee et controle d'acces par role;
  \item gestion des modules et de leurs composants pédagogiques (Cours/TD/TP);
  \item workflow de gestion des utilisateurs, assignations et inscriptions;
  \item publication/consultation des ressources pédagogiques;
  \item communication (chat) et suivi de progression;
  \item diffusion de sessions live avec acces cible aux apprenants concernes.
\end{itemize}

\subsection{Besoins non fonctionnels}

\begin{itemize}
  \item sécurité: JWT, validation des droits, isolation multi-école;
  \item performance: API REST reactive, requetes SQL optimisees;
  \item maintenabilite: architecture modulaire frontend/backend;
  \item evolutivite: ajout de nouvelles fonctionnalites sans rupture;
  \item ergonomie: réduction de la surcharge cognitive par role.
\end{itemize}

\subsection{Besoins detailles par acteur}

Afin de garantir une couverture fonctionnelle complete, les besoins ont ete declins acteur par acteur. Cette decomposition facilite la priorisation, la validation des droits et la preparation des tests de recette.

\renewcommand{\arraystretch}{1.35}
\begin{longtable}{|p{2.8cm}|p{4.6cm}|p{4.6cm}|}
\caption{Besoins detailles par acteur}\label{tab:besoins_acteurs}\\
\hline
\rowcolor{ustoblue}
\textcolor{white}{\textbf{Acteur}} &
\textcolor{white}{\textbf{Besoins prioritaires}} &
\textcolor{white}{\textbf{Valeur attendue}} \\
\hline
Super administrateur &
Creer les écoles, nommer les administrateurs, superviser la plateforme globale. &
Gouvernance transverse et maitrise des droits inter-etablissements. \\
\hline
Administrateur d'école &
Gérer comptes, modules, assignations, inscriptions, calendrier pédagogique. &
Pilotage quotidien de l’activité académique et réduction des taches manuelles. \\
\hline
Formateur Global &
Vérifie et valide les ajouts fait par Formateur Simple + même fonctionnalités que lui &
Production pédagogique structurée et interaction rapide avec les apprenants. \\
\hline
Formateur Simple &
Publier ressources, organiser Cours/TD/TP, lancer des sessions live, suivre les apprenants. &
Production pédagogique structurée et interaction rapide avec les apprenants. \\
\hline
Étudiant &
Consulter ressources, accéder aux lives, interagir via chat. &
Accès clair aux contenus et meilleure continuité d'apprentissage. \\
\hline
\end{longtable}

%%──────────────────────────────────────────────────────────────────────────────
\section{Modelisation UML}
%%──────────────────────────────────────────────────────────────────────────────

\subsection{Diagramme de cas d'utilisation}

Les cas d'utilisation centraux incluent:
\begin{itemize}
  \item \textbf{Administrer la plateforme}: creer écoles et administrateurs;
  \item \textbf{Administrer une école}: gerer utilisateurs, modules, assignations et inscriptions;
  \item \textbf{Produire et diffuser}: uploader et organiser les ressources;
  \item \textbf{Apprendre et interagir}: consulter, telecharger, completer des ressources;
  \item \textbf{Animer en direct}: creer, demarrer et terminer une session live.
\end{itemize}



\subsection{Cas d'utilisation detailles}

Pour augmenter la lisibilité et aligner la modélisation avec l'implementation réelle, le mémoire reserve trois diagrammes de cas d'utilisation distincts, chacun regroupant les roles fonctionnels suivants~:
\begin{itemize}
  \item administrateur (super administrateur et administrateur d'école);
  \item formateur (formateur global / principal et formateur simple);
  \item étudiant.
\end{itemize}

\subsubsection{Cas d'utilisation 1: Administrateur (super + école)}
\begin{figure}[H]
    \centering
    \safeimage[0.65\textwidth]{admin-usecase.png}{Cas d’utilisation: Administrateur}{Cas d’utilisation: Administrateur}
    \label{fig:usecase-admin}
\end{figure}
\clearpage

\subsubsection{Cas d'utilisation 2: Formateur}
\begin{figure}[H]
    \centering
    \safeimage[0.65\textwidth]{forma-usecase.png}{Cas d’utilisation: Formateur}{Cas d’utilisation: Formateur}
    \label{fig:usecase-formateur}
\end{figure}
\clearpage

\subsubsection{Cas d'utilisation 3: Etudiant}
\begin{figure}[H]
    \centering
    \safeimage[0.65\textwidth]{student-usecase.png}{Cas d'utilisation: Etudiant}{Cas d'utilisation: Etudiant}
    \label{fig:usecase-étudiant}
\end{figure}
\clearpage



\subsection{Diagrammes de séquence}

\subsubsection{Sequence 1: authentification et redirection par role}

\textbf{Participants:} Etudiant, Admin Ecole, Super Admin, Formateur, Syst\`eme EduPlatform.

\textbf{Flux nominal (code réel):}
\begin{enumerate}
  \item L'utilisateur accede \`a la plateforme.
  \item L'utilisateur saisit ses identifiants.
  \item Le backend verifie les informations et authentifie l'utilisateur.
  \item Le système identifie le role et cree la session.
  \item Le système redirige vers l'espace correspondant: étudiant, administration, formateur ou super admin.
\end{enumerate}

\begin{figure}[H]
    \centering
    \safeimage[0.72\textwidth]{seq1.png}{Diagramme de sequence: Authentification et redirection}{Diagramme de sequence: Authentification et redirection}
    \label{fig:seq1}
\end{figure}

\subsubsection{Sequence 2: creation de module et affectation d'un formateur}

\textbf{Participants:} Super Admin, Syst\`eme EduPlatform.

\textbf{Flux nominal (code réel):}
\begin{enumerate}
  \item Le super admin ouvre la gestion des modules.
  \item Il saisit les informations du nouveau module.
  \item Le système verifie les données et cree le module.
  \item Le super admin affecte un formateur au module.
  \item Le système verifie la disponibilite du formateur et associe le formateur au module.
  \item Le système confirme l'affectation reussie.
\end{enumerate}

\begin{figure}[H]
    \centering
    \safeimage[0.72\textwidth]{seq2.png}{Diagramme de sequence: Creation de module}{Diagramme de sequence: Creation de module}
    \label{fig:seq2}
\end{figure}
\clearpage

\subsubsection{Sequence 3: circulation de la demande d'inscription étudiante}

\textbf{Participants:} Etudiant, Admin Ecole, Syst\`eme EduPlatform.

\textbf{Flux nominal (code réel):}
\begin{enumerate}
  \item L'étudiant demande une inscription et envoie ses informations personnelles.
  \item Le système verifie le formulaire et cree une demande d'inscription.
  \item L'admin école consulte la demande et affiche les informations étudiant.
  \item Si la demande est acceptee, l'admin valide l'inscription et le système cree le compte étudiant.
  \item Si la demande est refusee, l'admin envoie la notification de refus.
\end{enumerate}

\begin{figure}[H]
    \centering
    \safeimage[0.72\textwidth]{seq3.png}{Diagramme de sequence: Demande d'inscription}{Diagramme de sequence: Demande d'inscription}
    \label{fig:seq3}
\end{figure}
\clearpage

\subsubsection{Sequence 4: publication de ressource pédagogique}

\textbf{Participants:} Formateur, Syst\`eme EduPlatform.

\textbf{Flux nominal (code réel):}
\begin{enumerate}
  \item Le formateur accede \`a ses modules affects.
  \item Il ajoute une ressource en envoyant le fichier et les informations associées.
  \item Le système verifie le fichier et les metadonnees.
  \item Si le fichier est valide, le système enregistre la ressource.
  \item Le formateur est informe de la reussite ou de l'erreur.
\end{enumerate}

\begin{figure}[H]
    \centering
    \safeimage[0.72\textwidth]{seq4.png}{Diagramme de sequence: Publication de ressource}{Diagramme de sequence: Publication de ressource}
    \label{fig:seq4}
\end{figure}
\clearpage

\subsubsection{Sequence 5: creation et cycle d'une session live}

\textbf{Participants:} Formateur, Etudiant, Syst\`eme EduPlatform.

\textbf{Flux nominal (code réel):}
\begin{enumerate}
  \item Le formateur cree une session live.
  \item Le système enregistre la seance et programme la notification.
  \item Le formateur demarre la seance en direct.
  \item L'étudiant rejoint la session, accede au direct et participe a la discussion.
  \item Le formateur termine la seance et le système sauvegarde l'historique.
\end{enumerate}

\begin{figure}[H]
    \centering
    \safeimage[0.72\textwidth]{seq5.png}{Diagramme de sequence: Session live}{Diagramme de sequence: Session live}
    \label{fig:seq5}
\end{figure}
\clearpage

\subsection{Diagramme de classes conceptuel}

Le modele conceptuel s'articule autour de classes principales:
\begin{itemize}
  \item \texttt{User} (roles, school, profil),
  \item \texttt{Module} (code, nom, niveau, composants),
  \item \texttt{Assignment} (formateur-module-type),
  \item \texttt{Enrollment} (étudiant-module-statut),
  \item \texttt{Resource} (categorie, fichier, statut),
  \item \texttt{LiveSession} (lien, statut, horodatages).
\end{itemize}

\begin{figure}[H]
    \centering
    \safeimage[0.65\textwidth]{umlclasse.drawio.png}{Diagramme de classes EduPlatform}{Diagramme de classes EduPlatform}
    \label{fig:view-adminmodule}
\end{figure}

%%──────────────────────────────────────────────────────────────────────────────
\section{Scenarios métier et captures d’écran}
%%──────────────────────────────────────────────────────────────────────────────

Cette section formalise les parcours métier clés (A à E) : création et gouvernance des modules, gestion des inscriptions, intégration d'assistants IA, flux d'approbation par le responsable de formation, et sessions live complètes.

\subsection{Scenario A: création d'un module et affectation d'un formateur}

\textbf{Préconditions:} l'administrateur de l'école est authentifié et dispose des droits de création de module.

\textbf{Flux principal:}
\begin{enumerate}
  \item l'administrateur saisit le formulaire de création (code, intitulé, niveau, composants, crédits);
  \begin{figure}[H]
    \centering
    \safeimage[0.65\textwidth]{Capture d’écran (307).png}{Vue Administrateur d'Ecole: Ajout de module}{Vue Administrateur d'Ecole: Ajout de module}
    \label{fig:view-adminmodule}
\end{figure}
  \item le système vérifie l'unicité du code au sein de l'école et normalise les composants;
  \item le module est enregistré dans la table \texttt{modules} (ou \texttt{module\_catalog} selon le déploiement);
  \begin{figure}[H]
    \centering
    \safeimage[0.65\textwidth]{Capture d’écran (309).png}{Vue Administrateur d'Ecole: Affectation de formateur}{Vue Administrateur d'Ecole: Affectation de formateur}
    \label{fig:view-adminmodule}
\end{figure}
  \item l'administrateur affecte un formateur principal (chargé de formation) et éventuellement des formateurs simples;
  \begin{figure}[H]
    \centering
    \safeimage[0.65\textwidth]{Capture d’écran (310).png}{Vue Formateur: Modules}{Vue Formateur: Modules}
    \label{fig:view-adminmodule}
\end{figure}
  \item le système envoie des notifications et crée la communauté/slug associée au module.
  \item le module devient visible aux acteurs autorisés (formateurs, étudiants inscrits).
\end{enumerate}

\textbf{Postconditions:} le module est actif, indexé et prêt à recevoir des ressources et des sessions.

\subsection{Scenario B: Le module chez le formateur}

\textbf{Préconditions:} le module est créé et un formateur (principal ou simple) est assigné au module.

\textbf{Flux principal:}
\begin{enumerate}
  \item le formateur consulte la fiche module depuis son tableau de bord;
  \item il accède aux composants du module (cours, TD, évaluations) et organise leur ordre et métadonnées;
  \item si le formateur est simple, il voit uniquement le \texttt{component\_scope} qui lui est assigné et ne peut modifier que ces composants;
  \begin{figure}[H]
    \centering
    \safeimage[0.9\textwidth]{Capture d’écran (323).png}{Vue Formateur: Modules}{Vue Formateur: Modules}
    \label{fig:view-adminmodule}
\end{figure}
\begin{figure}[H]
    \centering
    \safeimage[0.9\textwidth]{Capture d’écran (324).png}{Vue Formateur: Ajout ressource}{Vue Formateur: Ajout ressource}
    \label{fig:view-adminmodule}
\end{figure}
\end{enumerate}

\textbf{Postconditions:} le module est disponible et correctement structuré pour les contributions pédagogiques.

\subsection{Scenario C: Inscription d'un étudiant}

\textbf{Préconditions:} l'étudiant possède un compte valide; le module est disponible dans l'école.

\textbf{Flux principal:}
\begin{enumerate}
  \item l'étudiant (ou l'administrateur) remplit le formulaire d'inscription au module;
  \begin{figure}[H]
    \centering
    \safeimage[0.9\textwidth]{Capture d’écran (312).png}{Formulaire d'inscription}{Formulaire d'inscription}
    \label{fig:view-inscription-form}
\end{figure}
  \item la demande d'inscription est envoyée et traitée par l'administrateur de l'école;
  \begin{figure}[H]
    \centering
    \safeimage[0.9\textwidth]{Capture d’écran (313).png}{Vue Admin d'Ecole: Inscriptions}{Vue Admin d'Ecole: Inscriptions}
    \label{fig:view-admin-inscriptions}
\end{figure}
  \item le système vérifie les contraintes (même école, pas de doublon) et crée l'enregistrement dans \texttt{etudiant\_module\_enrollment};
  \item en cas de succès, l'étudiant voit le module sur son tableau de bord et reçoit les notifications.
\end{enumerate}

\textbf{Postconditions:} l'étudiant a accès aux ressources approuvées et reçoit les annonces et événements liés au module.

\subsection{Scenario D: communication et assistance via intelligence artificielle}

\textbf{Préconditions:} l'utilisateur (étudiant ou formateur) est authentifié; l'intégration IA (API interne ou service tiers) est disponible.

\textbf{Flux principal:}
\begin{enumerate}
  \item l'utilisateur ouvre le panneau d'assistance IA depuis une ressource ou la page module;
  \item il saisit une requête (explication d'un concept, génération de quiz, reformulation d'une fiche de cours);
  \begin{figure}[H]
    \centering
    \safeimage[0.65\textwidth]{Capture d’écran (325).png}{Vue Etudiant: Assistant IA}{Vue Etudiant: Assistant IA}
    \label{fig:view-etudiant-accueil}
\end{figure}
  \item la requête est envoyée au service IA, qui renvoie une réponse structurée (texte, questions, suggestions de ressources);
  \item le formateur peut accepter, éditer ou rejeter les suggestions; en cas d'acceptation il peut convertir la sortie IA en ressource ou activité publiable;
  \item les interactions avec l'IA sont journalisées pour traçabilité et contrôle qualité.
\end{enumerate}

\textbf{Postconditions:} suggestion IA enregistrée comme brouillon ou ressource (selon validation), avec historique d'audit.

% La sous-section "Scenario D" a été supprimée sur demande utilisateur.

\subsection{Scenario E: session live de bout en bout}

\textbf{Préconditions:} le formateur est assigné au module et dispose d'un lien de réunion valide; les étudiants sont informés de la session.

\textbf{Flux principal:}
\begin{enumerate}
  \item le formateur programme une session (statut \texttt{SCHEDULED}) avec date/heure et lien;
  \begin{figure}[H]
    \centering
    \safeimage[0.65\textwidth]{Capture d’écran (321).png}{Vue Formateur: Live}{Vue Formateur: Live}
    \label{fig:view-formateur-live}
\end{figure}
  \item une notification est envoyée aux étudiants inscrits et la session apparaît sur le dashboard;
  \item au moment du démarrage le formateur passe la session en statut \texttt{LIVE};
  \begin{figure}[H]
    \centering
    \safeimage[0.65\textwidth]{Capture d’écran (322).png}{Vue Etudiant: Acceuil}{Vue Etudiant: Acceuil}
    \label{fig:view-etudiant-accueil}
\end{figure}
  \item pendant la session, échanges synchrones (chat, Q\&A) et partages de ressources sont possibles; des marqueurs de présence/compréhension peuvent être collectés;
  \item à la fin, le responsable clôt la session en statut \texttt{ENDED} avec horodatage; le replay et les ressources associées sont archivés.
\end{enumerate}

\textbf{Postconditions:} session historisée, indicateurs d'engagement disponibles pour reporting pédagogique.

%%──────────────────────────────────────────────────────────────────────────────
\section{Conception des données}
%%──────────────────────────────────────────────────────────────────────────────

\subsection{Principes de conception}

Le schema relationnel repose sur:
\begin{itemize}
  \item normalisation des entites pour réduire la redondance;
  \item contraintes d'integrite (PK/FK, \texttt{CHECK}, \texttt{NOT NULL});
  \item journalisation temporelle (\texttt{created\_at}, \texttt{started\_at}, \texttt{ended\_at});
  \item indexation des colonnes critiques (recherche et tri par statut/date).
\end{itemize}

\subsection{Extension live\_sessions}

L'ajout de la table \texttt{live\_sessions} formalise la fonctionnalite de live pédagogique.
Les attributs cle incluent:
\begin{itemize}
  \item \texttt{formateur\_id}, \texttt{module\_id}, \texttt{meeting\_link};
  \item \texttt{status} dans \{SCHEDULED, LIVE, ENDED\};
  \item \texttt{scheduled\_at}, \texttt{started\_at}, \texttt{ended\_at}.
\end{itemize}

Ce schema garantit la traçabilite des sessions et facilite le filtrage par role.

\subsection{Regles de gestion et contraintes de données}

Pour renforcer la fiabilite metier, les regles suivantes ont ete retenues:
\begin{itemize}
  \item un module appartient a une seule école et son code y est unique;
  \item un étudiant ne peut pas etre inscrit deux fois au meme module;
  \item seuls les formateurs assignes peuvent creer ou modifier des ressources du module;
  \item une session live ne peut passer a \texttt{LIVE} que depuis l'etat \texttt{SCHEDULED};
  \item l'etat \texttt{ENDED} est terminal et interdit toute reprise de session;
  \item toute action d'administration doit etre rattachee a un utilisateur authentifie.
\end{itemize}

%%──────────────────────────────────────────────────────────────────────────────
\section{Architecture logicielle retenue}
%%──────────────────────────────────────────────────────────────────────────────

\subsection{Vue d'ensemble}

L'architecture suit un modele client-serveur decouple:
\begin{itemize}
  \item \textbf{Frontend}: React (SPA), routage protege, composants par role;
  \item \textbf{Backend}: Node.js/Express, controleurs REST, services metier, middlewares;
  \item \textbf{Base de données}: PostgreSQL, migrations SQL versionnees.
\end{itemize}

\subsection{Justification des choix}

\begin{itemize}
  \item React facilite une UX dynamique et modulaire;
  \item Express accelere la construction d'API REST maintenables;
  \item PostgreSQL assure robustesse transactionnelle et performance relationnelle;
  \item JWT permet une authentification stateless adaptee a une architecture API.
\end{itemize}

\begin{figure}[H]
    \centering
    \safeimage[0.9\textwidth]{archiEdu.drawio.png}{architecture-globale}{Architecture générale}
    \label{fig:architecture-globale}
\end{figure}

%%──────────────────────────────────────────────────────────────────────────────
\section{Risque projet et criteres d'acceptation}
%%──────────────────────────────────────────────────────────────────────────────

\subsection{Risques principaux}

\begin{itemize}
  \item complexite des droits multi-roles;
  \item derive UX en cas de surcharge d'ecrans admin;
  \item dette technique si l'évolution n'est pas guidee par conventions.
\end{itemize}

\subsection{Criteres d'acceptation}

\begin{itemize}
  \item chaque role accede uniquement aux fonctionnalites autorisees;
  \item les workflows critiques (module, inscription, live) sont executables de bout en bout;
  \item l'interface reste lisible sur desktop et mobile;
  \item les données critiques sont tracables et coherentes en base.
\end{itemize}

%%──────────────────────────────────────────────────────────────────────────────
\section{Matrice de tracabilite besoins-cas-tests}
%%──────────────────────────────────────────────────────────────────────────────

La tracabilite relie explicitement chaque besoin a un cas d'utilisation et a un test de validation. Cette pratique reduit le risque d'oublis fonctionnels durant la recette.

\renewcommand{\arraystretch}{1.35}
\begin{longtable}{|p{2.2cm}|p{4.1cm}|p{3.5cm}|p{3.6cm}|}
\caption{Matrice de tracabilite fonctionnelle}\label{tab:traceabilite}\\
\hline
\rowcolor{ustoblue}
\textcolor{white}{\textbf{Exigence}} &
\textcolor{white}{\textbf{Description}} &
\textcolor{white}{\textbf{Cas UML associe}} &
\textcolor{white}{\textbf{Validation attendue}} \\
\hline
OS1 & Centraliser ressources pédagogiques & Produire et diffuser & Ressources accessibles par module/composant \\
\hline
OS2 & Authentification et autorisation par roles & Administrer la plateforme/école & Acces refuses hors privileges \\
\hline
OS3 & Communication integree & Apprendre et interagir & Publications/chat operationnels \\
\hline
OS4 & Notifications des evenements & Apprendre et interagir & Notification visible apres evenement \\
\hline
OS5 & Tracabilite des operations & Administrer la plateforme/école & Historique coherent en base \\
\hline
Live & Session synchrone encadree & Animer en direct & Etats SCHEDULED/LIVE/ENDED conformes \\
\hline
\end{longtable}

%%──────────────────────────────────────────────────────────────────────────────
\section{Conclusion}
%%──────────────────────────────────────────────────────────────────────────────

Ce chapitre a etabli le pont entre besoins metier et solution technique. L'analyse fonctionnelle, les modeles UML et la conception de données offrent un cadre methodologique solide pour la realisation. Le chapitre suivant detaille l'implementation concrete, les choix de deploiement et les resultats de validation obtenus sur la plateforme.

%% ═══════════════════════════════════════════════════════════════════════════════
%% CHAPITRE 3 — Realisation, deploiement et validation
%% ═══════════════════════════════════════════════════════════════════════════════

\chapter{Implementation technique, deploiement et validation}

%%──────────────────────────────────────────────────────────────────────────────
\section{Introduction du chapitre}
%%──────────────────────────────────────────────────────────────────────────────

Apres la phase d'analyse et de conception, ce chapitre presente la realisation effective d'EduPlatform. Il decrit l'environnement de développement, l'implementation des modules fonctionnels, l'intégration de la fonctionnalite live, puis les tests de validation.

%%──────────────────────────────────────────────────────────────────────────────
\section{Environnement technologique}
%%──────────────────────────────────────────────────────────────────────────────

\subsection{Stack technique}

\begin{itemize}
  \item Frontend: React (SPA), React Router, Axios;
  \item Backend: Node.js, Express, middlewares de sécurité;
  \item Base de données: PostgreSQL;
  \item Authentification: JWT;
  \item IA locale: Ollama pour l'exécution de modèles de langage open source, la génération de contenus pédagogiques et l'analyse de ressources sans dépendance aux API cloud;
  \item Outils: VS Code, Git, scripts de migration SQL.
\end{itemize}

Ollama est utilisé dans le cadre d'un service IA embarqué qui peut exploiter des modèles open source en local ou sur une infrastructure contrôlée. Cette intégration permet de proposer plusieurs usages pédagogiques:
\begin{itemize}
  \item assistance à la rédaction et à la reformulation des fiches pédagogiques;
  \item génération d'exemples, de quiz ou de résumés à partir des ressources existantes;
  \item analyse des métadonnées et classification des contenus pour aider à la recherche par module et niveau;
  \item préservation de la confidentialité des données en limitant les échanges vers des services externes.
\end{itemize}

\subsection{Organisation du code}

Le projet est structure en deux blocs:
\begin{itemize}
  \item \textbf{frontend}: pages par role, composants UI, services API;
  \item \textbf{backend}: routes, controleurs, modeles, services metier.
\end{itemize}

Cette separation facilite la maintenabilite et l'évolution incrémentale.

\subsection{Configuration d'execution}

Le fonctionnement d'EduPlatform repose sur une separation claire des environnements de développement et de production.
\begin{itemize}
  \item \textbf{Variables d'environnement}: configuration des ports, URL API, secrets JWT et acces base de données;
  \item \textbf{Parametres de sécurité}: expiration des jetons, politique CORS et validation des entrees;
  \item \textbf{Scripts de lancement}: demarrage independant frontend/backend pour faciliter le debug par couche.
\end{itemize}

Cette configuration limite les regressions entre postes de développement et prepare la transition vers un deploiement reproductible.

%%──────────────────────────────────────────────────────────────────────────────
\section{Architecture d'execution et flux applicatifs}
%%──────────────────────────────────────────────────────────────────────────────

\subsection{Flux client-serveur principal}

Le cycle d'une requete suit les etapes suivantes:
\begin{enumerate}
  \item emission d'une action utilisateur depuis l'interface React;
  \item appel HTTP via service API avec jeton JWT si route protegee;
  \item passage par les middlewares de controle (authentification/autorisation/validation);
  \item execution de la logique metier dans les controleurs et services;
  \item acces PostgreSQL et retour JSON structure vers le frontend.
\end{enumerate}

\subsection{Gestion des erreurs}

Une strategie de gestion d'erreurs a ete appliquee pour garantir un comportement previsible:
\begin{itemize}
  \item codes HTTP coherents (401, 403, 404, 422, 500);
  \item messages explicites côté interface pour guider l'utilisateur;
  \item journalisation des erreurs serveur afin de faciliter le diagnostic.
\end{itemize}

%%──────────────────────────────────────────────────────────────────────────────
\section{Implementation des fonctionnalites principales}
%%──────────────────────────────────────────────────────────────────────────────

\subsection{Gestion des utilisateurs et des roles}

Le système applique un controle d'acces strict par role global et contexte école. Les routes critiques sont protegees par authentification et verification de privileges.

\subsection{Gestion pédagogique par module}

La plateforme gere:
\begin{itemize}
  \item creation/modification des modules;
  \item structuration par composants (Cours/TD/TP);
  \item assignation des formateurs (principal/simple);
  \item inscriptions des apprenants et suivi des ressources.
\end{itemize}

\subsection{Communication et collaboration}

Les utilisateurs disposent d'espaces de communication (chat/publications) afin de fluidifier les interactions pédagogiques et l'entraide.

\subsection{Traçabilite des operations metier}

Les operations critiques (creation module, assignation, inscription, changements d'etat live) sont realisees via des points d'entree explicites côté API. Cette structuration permet:
\begin{itemize}
  \item de verifier les droits avant toute modification en base;
  \item de conserver une chronologie claire des changements;
  \item de faciliter les audits fonctionnels en phase de validation.
\end{itemize}

%%──────────────────────────────────────────────────────────────────────────────
\section{Implementation de la fonctionnalite Live}
%%──────────────────────────────────────────────────────────────────────────────

\subsection{Conception metier du live}

La fonctionnalite live repose sur un cycle de vie simple:
\begin{itemize}
  \item \texttt{SCHEDULED}: session planifiee;
  \item \texttt{LIVE}: session en direct;
  \item \texttt{ENDED}: session terminee.
\end{itemize}

Le formateur cree une session avec un lien de reunion (Zoom/Meet/Teams), puis controle son etat. Les apprenants concernes voient automatiquement les sessions actives ou programmees.

\subsection{Realisation backend}

Les points d'API implementes couvrent:
\begin{itemize}
  \item creation et consultation des sessions du formateur;
  \item transition d'etat (demarrer/terminer);
  \item suppression d'une session;
  \item recuperation des lives pertinents côté étudiant.
\end{itemize}

Pour renforcer la robustesse, le backend verifie egalement la cohérence des transitions d'etat. Une session ne peut pas passer directement de \texttt{SCHEDULED} a \texttt{ENDED} sans passage par \texttt{LIVE}, sauf action d'annulation explicite.

\subsection{Realisation frontend}

Deux integrations ont ete realisees:
\begin{itemize}
  \item page dediee formateur pour gerer les sessions live;
  \item widget dashboard étudiant affichant les lives actifs/programmes.
\end{itemize}

Cette implementation assure une liaison directe entre la planification enseignant et l'experience apprenant.

%%──────────────────────────────────────────────────────────────────────────────
\section{Securite applicative et protection des données}
%%──────────────────────────────────────────────────────────────────────────────

La sécurité a ete traitee comme une contrainte transversale, non comme une fonctionnalite secondaire.

\subsection{Authentification et autorisation}

\begin{itemize}
  \item authentification par jeton JWT signe;
  \item verification des droits selon role global et contexte école;
  \item protection des routes d'administration et d'ecriture.
\end{itemize}

\subsection{Validation des données et réduction des risques}

\begin{itemize}
  \item validation serveur des champs obligatoires et formats attendus;
  \item controles metier sur les relations formateur-module-étudiant;
  \item limitation des effets de bord grace a des flux CRUD explicites.
\end{itemize}

\subsection{Protection operationnelle minimale}

Les mesures suivantes ont ete appliquees dans le cadre du projet:
\begin{itemize}
  \item restriction des operations sensibles aux profils autorises;
  \item gestion de l'expiration des sessions d'authentification;
  \item usage de configurations separees pour eviter l'exposition de secrets en code source.
\end{itemize}

%%──────────────────────────────────────────────────────────────────────────────
\section{Deploiement et exploitation}
%%──────────────────────────────────────────────────────────────────────────────

\subsection{Strategie de deploiement retenue}

Le deploiement est organise en trois niveaux:
\begin{itemize}
  \item \textbf{Developpement local}: iteration rapide des fonctionnalites;
  \item \textbf{Pre-production}: verification des parcours critiques avec données de test;
  \item \textbf{Production pilote}: ouverture controlee aux utilisateurs reels.
\end{itemize}

\subsection{Sauvegarde et continuité}

Afin de limiter le risque de perte de données:
\begin{itemize}
  \item des sauvegardes periodiques de la base PostgreSQL sont planifiees;
  \item les scripts de restauration sont documentes pour une reprise rapide;
  \item les evolutions schema sont tracees via scripts SQL versionnes.
\end{itemize}

\subsection{Supervision et maintenance}

La maintenance applicative suit une boucle continue:
\begin{enumerate}
  \item collecte des incidents et retours utilisateurs;
  \item priorisation selon impact pédagogique et frequence;
  \item correction, verification et redeploiement cible.
\end{enumerate}

%%──────────────────────────────────────────────────────────────────────────────
\section{Ameliorations UX realisees}
%%──────────────────────────────────────────────────────────────────────────────

Une passe UX a ete effectuee pour réduire la surcharge cognitive:
\begin{itemize}
  \item affichage des modules admin en grille lisible;
  \item filtres metier (niveau, role, module, type) sur ecrans admin;
  \item suppression de colonnes peu utiles (exemple: ID) dans les tableaux;
  \item enrichissement visuel des tableaux de bord et cartes de suivi;
  \item intégration visuelle du bloc live sur les dashboards concernes.
\end{itemize}

%%──────────────────────────────────────────────────────────────────────────────
\section{Validation fonctionnelle}
%%──────────────────────────────────────────────────────────────────────────────

\subsection{Jeux d'essais}

Les validations ont ete conduites sur des comptes de test multi-roles:
\begin{itemize}
  \item administration (super-admin et admin école),
  \item formateurs,
  \item apprenants.
\end{itemize}

\subsection{Resultats observes}

\begin{itemize}
  \item authentification et redirections par role conformes;
  \item workflows de gestion (module/assignation/inscription) operationnels;
  \item fonctionnalite live testee de bout en bout (creation -> LIVE -> affichage étudiant);
  \item compilation frontend validee sans erreurs bloquantes sur les ecrans modifies.
\end{itemize}

\subsection{Tableau synthetique de validation}

\renewcommand{\arraystretch}{1.3}
\begin{table}[H]
\centering
\caption{Synthese des tests fonctionnels principaux}
\label{tab:tests_ch3}
\begin{tabular}{|p{4.3cm}|p{3.2cm}|p{2.4cm}|p{2.6cm}|}
\hline
\rowcolor{ustoblue}
\textcolor{white}{Scenario teste} &
\textcolor{white}{Role principal} &
\textcolor{white}{Statut} &
\textcolor{white}{Observation} \\
\hline
Creation/edition module & Admin école & Valide & Workflow stable \\
\hline
Assignation formateur & Admin école & Valide & Regles de droits respectees \\
\hline
Inscription étudiant & Admin école & Valide & Visibilite module immediate \\
\hline
Publication ressource & Formateur & Valide & Organisation Cours/TD/TP conforme \\
\hline
Session live complete & Formateur + Etudiant & Valide & Etats SCHEDULED/LIVE/ ENDED coherents \\
\hline
Consultation ressources & Etudiant & Valide & Parcours fluide et lisible \\
\hline
\end{tabular}
\end{table}

\subsection{Limites actuelles}

\begin{itemize}
  \item les liens live sont externes (pas de visio embarquee native);
  \item la metrologie produit (analytics avances) reste a industrialiser;
  \item certains tests automatises restent a renforcer pour le passage a l'echelle.
\end{itemize}

\subsection{Axes d'amelioration technique court terme}

\begin{itemize}
  \item augmenter la couverture de tests automatises sur les routes critiques;
  \item introduire des indicateurs de performance (latence API, erreurs par endpoint);
  \item renforcer la documentation d'exploitation pour faciliter l'onboarding technique.
\end{itemize}

%%──────────────────────────────────────────────────────────────────────────────
\section{Synthese et perspectives techniques}
%%──────────────────────────────────────────────────────────────────────────────

La realisation confirme la faisabilite technique du socle EduPlatform et sa capacite a supporter des usages reels multi-acteurs. Les evolutions prioritaires envisagees sont:
\begin{itemize}
  \item observabilite avancee (logs metier, tableaux de supervision),
  \item automatisation QA (tests intégration/e2e),
  \item optimisation des performances sur charge elevee.
\end{itemize}

%%──────────────────────────────────────────────────────────────────────────────
\section{Conclusion}
%%──────────────────────────────────────────────────────────────────────────────

Ce chapitre a decrit l'implementation concrete d'EduPlatform et les validations associees. Les resultats montrent une plateforme fonctionnelle, évolutive et coherente avec les objectifs du projet. Le chapitre suivant complete cette analyse par la dimension entrepreneuriale: modele économique, SWOT et strategie de croissance startup.

%% ═══════════════════════════════════════════════════════════════════════════════
%% CHAPITRE 4 — Modele économique, positionnement startup et analyse SWOT
%% ═══════════════════════════════════════════════════════════════════════════════

\chapter{Evaluation stratégique et business model (cadre startup, loi 12-75)}

%%──────────────────────────────────────────────────────────────────────────────
\section{Introduction du chapitre}
%%──────────────────────────────────────────────────────────────────────────────

Dans les chapitres precedents, nous avons expose le contexte, la problématique, puis les choix techniques et fonctionnels qui structurent la plateforme EduPlatform. Le present chapitre deplace l'analyse vers la viabilite entrepreneuriale du projet. L'objectif est de verifier que la proposition de valeur n'est pas uniquement techniquement realisable, mais egalement economiquement soutenable dans un contexte de startup EdTech.

Ce chapitre poursuit quatre finalites. Premierement, formaliser un modele économique coherent a travers le \textit{Business Model Canvas} (BMC)\cite{osterwalder2010bmc}. Deuxiemement, evaluer la position stratégique de la startup au moyen d'une matrice SWOT\cite{kotler2017marketing}. Troisiemement, proposer des scenarios de monetisation et des indicateurs de pilotage. Quatriemement, definir une feuille de route de croissance progressive adaptee au contexte des écoles et centres de formation.

%%──────────────────────────────────────────────────────────────────────────────
\section{Cadre juridique et regulatoire}

Dans le cadre national, le projet doit prendre en compte le contexte reglémentaire applicable aux services numeriques et aux startups EdTech. La loi \textbf{12-75} figure parmi les textes de reference qui structurent l'environnement juridique des systèmes d'information et des activites economiques en lien avec les technologies de l'information. EduPlatform est ainsi confronte a des exigences de conformité sur les plans de la protection des données, de la sécurité des infrastructures et de la contractualisation des services.

Ce cadre juridique renforce l'importance d'un modele économique clairement defini, d'une gouvernance documentee et d'une roadmap de deploiement qui integre les aspects reglementaires des accords avec les etablissements.

%%──────────────────────────────────────────────────────────────────────────────
\section{Vision startup et hypotheses de depart}
%%──────────────────────────────────────────────────────────────────────────────

EduPlatform est envisage comme une startup B2B2C EdTech. Le coeur du modele consiste a contractualiser avec les etablissements (B2B), tout en delivrant de la valeur directe aux formateurs et apprenants (B2C indirect). Ce positionnement permet de traiter simultanement les enjeux institutionnels (gouvernance, sécurité, traçabilite) et les usages quotidiens (cours, TD, TP, lives, communication, suivi).

\subsection{Probleme marche adresse}

Les écoles de soutien, les écoles de langues et les établissements privés font face à trois contraintes récurrentes:
\begin{itemize}
  \item dispersion des ressources pédagogiques entre plusieurs canaux non gouvernés;
  \item manque de pilotage centralisé (droits, validations, historique, reporting);
  \item faible qualité de l'expérience numérique pour les enseignants et apprenants.
\end{itemize}

\subsection{Hypotheses business initiales}

\begin{itemize}
  \item Le besoin de centralisation pédagogique est suffisamment critique pour justifier un abonnement institutionnel.
  \item Les etablissements privilegient une solution locale, configurable et conforme aux exigences de gouvernance.
  \item Les fonctionnalites differenciantes (multi-école, workflow de validation, sessions live, suivi) augmentent le taux d'adoption interne.
  \item Un deploiement progressif par pilote (une école, puis extension) reduit le risque commercial et technique.
\end{itemize}

\subsection{Hypotheses de validation produit-marche}

Au-dela des hypotheses business, trois hypotheses de validation produit-marche sont retenues:
\begin{itemize}
  \item \textbf{H1 --- Adoption formateur}: au moins 70\,\% des formateurs pilotes publient au minimum une ressource par semaine.
  \item \textbf{H2 --- Activation apprenant}: au moins 60\,\% des apprenants connectes consultent activement les ressources sur une periode de 30 jours.
  \item \textbf{H3 --- Valeur institutionnelle}: l'administration confirme une réduction perceptible du temps de coordination pédagogique.
\end{itemize}

Ces hypotheses servent de base a la priorisation roadmap et a la décision de passage a l'echelle.

%%──────────────────────────────────────────────────────────────────────────────
\section{Analyse du marche et positionnement concurrentiel}
%%──────────────────────────────────────────────────────────────────────────────

Le positionnement d'EduPlatform ne se limite pas a une comparaison fonctionnelle avec les LMS etablis. Il repose sur une strategie de specialisation locale, orientee gouvernance pédagogique et simplicite d'usage.

\subsection{Lecture TAM / SAM / SOM}

\begin{itemize}
  \item \textbf{TAM (Total Addressable Market)}: l'ensemble des etablissements d'enseignement superieur et structures de formation pouvant numeriser leurs processus pédagogiques.
  \item \textbf{SAM (Serviceable Available Market)}: les institutions disposant d'une maturite numerique minimale et d'un besoin explicite de plateforme unifiee.
  \item \textbf{SOM (Serviceable Obtainable Market)}: la part realiste atteignable a court terme via pilotes regionaux et effet de recommandation inter-etablissements.
\end{itemize}

Cette lecture permet d'eviter une surprojection commerciale et d'ancrer les objectifs de croissance dans une trajectoire executable.

\subsection{Differenciation concurrentielle}

Les concurrents generalistes couvrent un spectre large, mais avec une adaptation locale limitee. EduPlatform se differencie sur:
\begin{itemize}
  \item la gouvernance fine multi-école et multi-roles;
  \item la structuration pédagogique native Cours/TD/TP;
  \item la proximite terrain pour l'accompagnement et l'iteration produit;
  \item la capacite d'implementation progressive sans rupture organisationnelle.
\end{itemize}

\renewcommand{\arraystretch}{1.3}
\begin{table}[H]
\centering
\caption{Positionnement concurrentiel simplifie}
\label{tab:positionnement_concurrentiel}
\begin{tabular}{|p{3.0cm}|p{3.2cm}|p{3.2cm}|p{3.2cm}|}
\hline
\rowcolor{ustoblue}
\textcolor{white}{Solution} & \textcolor{white}{Adaptation locale} & \textcolor{white}{Gouvernance pédagogique} & \textcolor{white}{Accompagnement} \\
\hline
LMS internationaux & Moyenne a faible & Bonne mais generaliste & Standardise \\
\hline
Outils bureautiques + chat & Faible & Fragmentee & Informel \\
\hline
EduPlatform & Elevee & Native au contexte metier & Proximite + iteration \\
\hline
\end{tabular}
\end{table}

%%──────────────────────────────────────────────────────────────────────────────
\section{Business Model Canvas d'EduPlatform}
%%──────────────────────────────────────────────────────────────────────────────

Le BMC formalise les neuf blocs strategiques de la startup\cite{osterwalder2010bmc}.

\subsection{Segments de clients}

EduPlatform cible quatre segments prioritaires:
\begin{itemize}
  \item écoles, centres de formation et etablissements de soutien cherchant une plateforme pédagogique gouvernee;
  \item administrations pédagogiques (direction, scolarite, responsables filiere);
  \item corps enseignant (formateurs principaux et simples);
  \item apprenants, utilisateurs finaux de l'expérience d'apprentissage.
\end{itemize}

\subsection{Proposition de valeur}

La proposition de valeur s'articule autour de six axes:
\begin{itemize}
  \item centralisation des contenus par niveau, module et composant (Cours/TD/TP);
  \item gouvernance fine des acces par roles et écoles;
  \item interactions pédagogiques integrees (chat, ressources, lives, notifications);
  \item réduction de la surcharge cognitive via une interface structuree;
  \item visibilité admin avec tableaux de bord et actions rapides;
  \item architecture évolutive orientee startup (MVP, iterations, extension multi-tenant).
\end{itemize}

\subsection{Canaux}

\begin{itemize}
  \item canal direct: demonstrations, POC et pilotes dans les etablissements;
  \item canal relationnel: réseau éducatif local, enseignants relais, bouche-à-oreille;
  \item canal digital: site vitrine, documentation, supports video de prise en main;
  \item canal partenarial: incubateurs, associations EdTech, acteurs institutionnels.
\end{itemize}

\subsection{Relations clients}

\begin{itemize}
  \item accompagnement de lancement (onboarding admin et formateurs);
  \item support fonctionnel continu (FAQ, tickets, sessions de formation);
  \item co-construction produit (collecte feedback, roadmap trimestrielle);
  \item contrats annuels avec points de revues periodiques.
\end{itemize}

\subsection{Flux de revenus}

Le modele de revenus privilegie l'abonnement institutionnel annuel:
\begin{itemize}
  \item formule \textbf{Starter}: petite école, fonctions essentielles;
  \item formule \textbf{Growth}: multi-filieres, reporting avance, priorite support;
  \item formule \textbf{Enterprise}: forte volumetrie, SLA, personnalisation avancee.
\end{itemize}

Des revenus complementaires peuvent etre ajoutes: formation premium, migration de données, developpements specifiques, et support etendu.

\subsection{Ressources cles}

\begin{itemize}
  \item équipe produit (développement full-stack, QA, design UX);
  \item infrastructure cloud/serveur et base PostgreSQL;
  \item bibliotheque de composants et architecture API;
  \item savoir-faire pédagogique et connaissance du terrain éducatif.
\end{itemize}

\subsection{Activites cles}

\begin{itemize}
  \item développement et maintenance applicative;
  \item securisation et supervision de la plateforme;
  \item deploiement, parametrage et support client;
  \item analyse d'usage et priorisation roadmap.
\end{itemize}

\subsection{Partenaires cles}

\begin{itemize}
  \item etablissements pilotes (co-validation des usages);
  \item hebergeurs/infrastructure cloud;
  \item ecosysteme startup (incubateurs, mentors, experts juridiques);
  \item partenaires de formation et d'intégration.
\end{itemize}

\subsection{Structure de couts}

\begin{itemize}
  \item couts techniques: hebergement, sauvegarde, monitoring, sécurité;
  \item couts humains: développement, support, commercial, pilotage produit;
  \item couts operationnels: deplacements, documentation, formation clients;
  \item couts d'acquisition: demos, communication, prospection.
\end{itemize}

\subsection{Synthese du BMC}

\renewcommand{\arraystretch}{1.45}
\begin{longtable}{|>{\bfseries\raggedright\arraybackslash}p{3.2cm}|p{10.2cm}|}
\caption{Synthese du Business Model Canvas d'EduPlatform}
\label{tab:bmc_edconnect} \\
\hline
\rowcolor{ustoblue}
\textcolor{white}{Bloc} & \textcolor{white}{Elements principaux} \\
\hline
Segments clients & Écoles/centres de formation, administration pédagogique, formateurs, apprenants \\
\hline
Proposition de valeur & Plateforme pédagogique unifiee, gouvernance par roles, modules Cours/TD/TP, live et collaboration \\
\hline
Canaux & Pilotes institutionnels, réseau éducatif local, canal digital, partenariats \\
\hline
Relations clients & Onboarding, support continu, co-construction produit, revues periodiques \\
\hline
Revenus & Abonnements annuels (Starter, Growth, Enterprise), services premium \\
\hline
Ressources cles & Equipe tech/produit, infrastructure, architecture logicielle, expertise metier \\
\hline
Activites cles & Dev/maintenance, sécurité, support, deploiement, analyse d'usage \\
\hline
Partenaires cles & Ecoles pilotes, hebergeurs, ecosysteme startup, partenaires intégration \\
\hline
Structure de couts & Technique, équipe, operations, acquisition client \\
\hline
\end{longtable}

%%──────────────────────────────────────────────────────────────────────────────
\section{Analyse SWOT de la startup EduPlatform}
%%──────────────────────────────────────────────────────────────────────────────

L'analyse SWOT permet de croiser facteurs internes (forces, faiblesses) et externes (opportunites, menaces) pour orienter les décisions strategiques.

\renewcommand{\arraystretch}{1.45}
\begin{table}[H]
\centering
\caption{Matrice SWOT d'EduPlatform}
\label{tab:swot_edconnect}
\begin{tabular}{|p{6.2cm}|p{6.2cm}|}
\hline
\rowcolor{ustoblue}
\textcolor{white}{\textbf{Forces (S)}} & \textcolor{white}{\textbf{Faiblesses (W)}} \\
\hline
\begin{itemize}
  \item Architecture multi-école et gouvernance par roles.
  \item UX adaptee au contexte pédagogique local.
  \item Fonctionnalites differenciantes: live, workflow, separation Cours/TD/TP.
  \item Équipe proche du terrain éducatif local.
\end{itemize}
&
\begin{itemize}
  \item Notoriete de marque encore faible.
  \item Ressources humaines et budget limites en phase initiale.
  \item Dependance a quelques profils techniques cles.
  \item Processus commerciaux encore en structuration.
\end{itemize}
\\
\hline
\rowcolor{ustoblue}
\textcolor{white}{\textbf{Opportunites (O)}} & \textcolor{white}{\textbf{Menaces (T)}} \\
\hline
\begin{itemize}
  \item Acceleration de la digitalisation des etablissements.
  \item Besoin croissant d'outils de souverainete et de gouvernance des données.
  \item Programmes d'appui a l'innovation et incubateurs.
  \item Faible adequation des solutions generalistes aux contraintes locales.
\end{itemize}
&
\begin{itemize}
  \item Concurrence LMS internationaux etablis.
  \item Ralentissements administratifs dans la décision d'achat.
  \item Risques cyber et exigences de conformite.
  \item Volatilite budgetaire des institutions.
\end{itemize}
\\
\hline
\end{tabular}
\end{table}

\subsection{Strategies derivees de la SWOT}

\paragraph{Strategie SO (exploiter forces + opportunites).}
Prioriser les deploiements pilotes dans des etablissements ouverts a l'innovation et transformer les resultats obtenus en references commerciales.

\paragraph{Strategie ST (utiliser forces contre menaces).}
Renforcer les arguments de differenciation locale (gouvernance, adaptation pédagogique, support de proximite) face aux LMS internationaux standardises.

\paragraph{Strategie WO (réduire faiblesses via opportunites).}
Utiliser l'appui d'incubateurs et partenaires pour structurer la fonction commerciale, formaliser les offres et accelerer l'acquisition.

\paragraph{Strategie WT (limiter faiblesses et menaces).}
Mettre en place une feuille de route sécurité/conformite, une documentation d'exploitation et une politique de continuité afin de réduire les risques operationnels.

%%──────────────────────────────────────────────────────────────────────────────
\section{Scenario de monetisation et indicateurs de pilotage}
%%──────────────────────────────────────────────────────────────────────────────

\subsection{Scenario de monetisation recommande}

Le scenario recommande est un abonnement annuel institutionnel avec segmentation par taille et besoin. Cette approche simplifie la lisibilite tarifaire et facilite le pilotage de la marge.

\begin{table}[H]
\centering
\caption{Structure d'offre proposee (exemple de cadrage)}
\label{tab:offre_pricing}
\begin{tabular}{|p{3.2cm}|p{3.4cm}|p{4.0cm}|p{2.0cm}|}
\hline
\rowcolor{ustoblue}
\textcolor{white}{Formule} & \textcolor{white}{Cible} & \textcolor{white}{Valeur livree} & \textcolor{white}{Mode de prix} \\
\hline
Starter & Petite école / pilote & Modules, ressources, roles, support standard & Abonnement annuel \\
\hline
Growth & Ecole en expansion & Multi-filiere, reporting, priorite support & Abonnement annuel \\
\hline
Enterprise & Grand etablissement & SLA, personnalisation, accompagnement etendu & Contrat annuel negocie \\
\hline
\end{tabular}
\end{table}

\subsection{Indicateurs cle (KPI startup)}

\begin{itemize}
  \item \textbf{KPI acquisition}: nombre d'etablissements pilotes signes, taux de conversion demo $\rightarrow$ contrat.
  \item \textbf{KPI adoption}: taux d'activation des formateurs, taux d'activite étudiant, volume de ressources publiees.
  \item \textbf{KPI retention}: renouvellement annuel, churn institutionnel, satisfaction per role.
  \item \textbf{KPI produit}: disponibilite plateforme, temps moyen de resolution incident, usage des lives.
  \item \textbf{KPI financier}: revenu recurrent annuel (ARR), cout d'acquisition client (CAC), marge brute.
\end{itemize}

%%──────────────────────────────────────────────────────────────────────────────
\section{Feuille de route de deploiement startup}
%%──────────────────────────────────────────────────────────────────────────────

\subsection{Phase 1: Validation pilote (0--6 mois)}

\begin{itemize}
  \item selection de 1 a 2 etablissements pilotes;
  \item contractualisation simple et accompagnement fort;
  \item mesure des KPI d'adoption et ajustements UX prioritaires.
\end{itemize}

\subsection{Phase 2: Structuration commerciale (6--12 mois)}

\begin{itemize}
  \item formalisation des trois offres;
  \item creation des supports commerciaux (demo, fiche valeur, FAQ);
  \item extension a de nouveaux etablissements via references pilotes.
\end{itemize}

\subsection{Phase 3: Passage a l'echelle (12--24 mois)}

\begin{itemize}
  \item industrialisation de l'onboarding client;
  \item renforcement sécurité/conformite et reporting avance;
  \item croissance geographique et sectorielle progressive.
\end{itemize}

\subsection{Plan go-to-market operationnel}

Le deploiement commercial peut etre structure en quatre chantiers synchronises:
\begin{enumerate}
  \item \textbf{Ciblage et prospection}: cartographier les etablissements prioritaires et qualifier les decideurs.
  \item \textbf{Preuve de valeur}: realiser des demonstrations orientees cas d'usage et indicateurs mesurables.
  \item \textbf{Onboarding pilote}: accompagner intensivement les 60 premiers jours pour garantir l'adoption.
  \item \textbf{Conversion et extension}: transformer les pilotes reussis en contrats annuels et recommandations.
\end{enumerate}

%%──────────────────────────────────────────────────────────────────────────────
\section{Conclusion}
%%──────────────────────────────────────────────────────────────────────────────

L'analyse entrepreneuriale montre qu'EduPlatform dispose d'une base solide pour evoluer d'un projet académique vers une startup EdTech viable. Le Business Model Canvas met en evidence un positionnement clair B2B2C, tandis que la SWOT confirme une fenetre d'opportunite réelle, a condition de structurer rapidement les dimensions commerciales, organisationnelles et securitaires.

Ce chapitre etablit ainsi le lien entre la pertinence pédagogique de la solution et sa soutenabilite économique. Il constitue une etape essentielle avant la phase d'industrialisation et de mise sur le marche, en orientant les décisions de priorisation produit, d'acquisition client et de gouvernance de croissance.

%% ═══════════════════════════════════════════════════════════════════════════════
%% CONCLUSION GENERALE
%% ═══════════════════════════════════════════════════════════════════════════════

\chapter*{Conclusion générale}
\addcontentsline{toc}{chapter}{Conclusion générale}

Ce mémoire avait pour objectif de concevoir et realiser une plateforme pédagogique unifiee repondant aux besoins d'un contexte éducatif réel. La démarche adoptee a permis de couvrir l'ensemble de la chaine de valeur du projet: analyse du contexte, modelisation, implementation technique, validation fonctionnelle et projection startup.

Sur le plan technique, EduPlatform propose un socle coherent et operationnel: gestion multi-roles, administration pédagogique, structuration des contenus par composants (Cours/TD/TP), communication integree, suivi et fonctionnalite live. Les tests fonctionnels ont confirme la stabilite des parcours critiques pour les profils administrateur, formateur et étudiant.

Sur le plan methodologique, la structuration du mémoire s'aligne sur les references d'encadrement (branding et bibliotheque virtuelle), tout en introduisant une extension majeure: la dimension entrepreneuriale. Le chapitre dedie au modele économique et a l'analyse SWOT montre qu'EduPlatform peut evoluer d'un projet académique vers une startup EdTech viable, sous reserve d'une execution disciplinee sur les axes produit, commercial et operationnel.

\paragraph{Contributions principales.}
\begin{itemize}
  \item definition d'une architecture pédagogique unifiee adaptee au terrain local;
  \item implementation d'une plateforme multi-acteurs avec gouvernance par role;
  \item intégration d'une fonctionnalite live orientee usage réel;
  \item formalisation d'un cadre business startup (BMC, SWOT, KPI, roadmap).
\end{itemize}

\paragraph{Perspectives.}
Les perspectives prioritaires concernent:
\begin{itemize}
  \item le renforcement des tests automatisees et de l'observabilite;
  \item l'industrialisation de l'onboarding institutionnel;
  \item l'intégration progressive d'outils IA pédagogiques;
  \item l'extension commerciale vers d'autres etablissements.
\end{itemize}

En synthese, EduPlatform demontre qu'une solution locale, évolutive et orientee usages peut concilier rigueur académique, efficacite operationnelle et potentiel de croissance startup.


\clearpage
\printbibliography[title={Bibliographie}, heading=bibintoc]

\end{document}

